\documentclass[12pt,a4paper]{ctexart}

% 页面设置
\usepackage[
    left=2.5cm,
    right=2.5cm,
    top=2.3cm,
    bottom=2.3cm
]{geometry}

% 常用宏包
\usepackage{xcolor}
\usepackage{tcolorbox}
\usepackage{titlesec}
\usepackage{fancyhdr}
\usepackage{enumitem}
\usepackage{tabularx}
\usepackage{amsmath, amssymb}
\usepackage{hyperref}

\tcbuselibrary{skins, breakable}

% ======================
% 颜色主题：清爽学习笔记风
% ======================

% 主色：深蓝绿
\definecolor{MainColor}{HTML}{2A6F73}

% 辅助色：青绿色
\definecolor{SecondColor}{HTML}{3CAEA3}

% 浅背景色
\definecolor{LightMain}{HTML}{EAF7F6}

% 标题深色
\definecolor{TitleDark}{HTML}{1F3A40}

% 文字灰
\definecolor{TextGray}{HTML}{555555}

% 边框灰
\definecolor{BorderGray}{HTML}{CBD5D6}

% 提醒色：柔和橙
\definecolor{WarnColor}{HTML}{F2B35D}

% 易错点红色
\definecolor{ErrorColor}{HTML}{C94C4C}

% 定义框紫色
\definecolor{DefineColor}{HTML}{6C63B7}

% ======================
% 正文设置
% ======================

\linespread{1.25}
\setlength{\parindent}{2em}
\setlength{\parskip}{0.25em}

\hypersetup{
  colorlinks=true,
  linkcolor=MainColor,
  urlcolor=MainColor
}

% ======================
% 页眉页脚
% ======================

\pagestyle{fancy}
\fancyhf{}
\fancyfoot[L]{\textcolor{MainColor}{机器学习笔记}}
\fancyfoot[R]{\textcolor{MainColor}{\thepage}}
\renewcommand{\headrulewidth}{0pt}
\renewcommand{\footrulewidth}{0pt}

% ======================
% 顶部小标题栏
% ======================

\newcommand{\topbar}[2]{
  \noindent
  {\small\textcolor{MainColor}{#1}}
  \hfill
  {\small\bfseries\textcolor{TitleDark}{#2}}
  \par\vspace{0.35em}
  {\color{BorderGray}\hrule}
  \vspace{1.2em}
}

% ======================
% 标题卡片
% ======================

\newcommand{\notetitlecard}[6]{
\begin{tcolorbox}[
  enhanced,
  colback=white,
  colframe=BorderGray,
  boxrule=0.6pt,
  arc=2mm,
  left=12pt,
  right=12pt,
  top=10pt,
  bottom=10pt,
  borderline west={4pt}{0pt}{SecondColor}
]
\begin{tabularx}{\linewidth}{X r}
{\small\textcolor{TextGray}{#1}} & {\small\textcolor{TextGray}{#2}} \\[0.9em]
\multicolumn{2}{l}{
  {\Huge\bfseries\textcolor{MainColor}{#3}}
} \\[0.45em]
\multicolumn{2}{l}{
  {\large\textcolor{SecondColor!90!black}{#4}}
} \\[0.45em]
\multicolumn{2}{l}{
  {\small\textcolor{TextGray}{课程：#5 \qquad 作者：#6}}
}
\end{tabularx}
\end{tcolorbox}
\vspace{1em}
}

% ======================
% 摘要框
% ======================

\newtcolorbox{summarybox}[1][本节摘要]{
  enhanced,
  breakable,
  colback=LightMain,
  colframe=SecondColor,
  colbacktitle=SecondColor,
  coltitle=white,
  title={#1},
  fonttitle=\bfseries,
  boxrule=0.6pt,
  arc=2mm,
  left=12pt,
  right=12pt,
  top=10pt,
  bottom=10pt
}

% ======================
% 定义框
% ======================

\newtcolorbox{definitionbox}[1][定义]{
  enhanced,
  breakable,
  colback=DefineColor!6,
  colframe=DefineColor,
  colbacktitle=DefineColor,
  coltitle=white,
  title={#1},
  fonttitle=\bfseries,
  boxrule=0.6pt,
  arc=2mm,
  left=12pt,
  right=12pt,
  top=10pt,
  bottom=10pt
}

% ======================
% 重点框
% ======================

\newtcolorbox{keybox}[1][重点]{
  enhanced,
  breakable,
  colback=MainColor!5,
  colframe=MainColor,
  colbacktitle=MainColor,
  coltitle=white,
  title={#1},
  fonttitle=\bfseries,
  boxrule=0.6pt,
  arc=2mm,
  left=12pt,
  right=12pt,
  top=10pt,
  bottom=10pt
}

% ======================
% 例题框
% ======================

\newtcolorbox{examplebox}[1][例题]{
  enhanced,
  breakable,
  colback=white,
  colframe=SecondColor,
  colbacktitle=SecondColor!85!black,
  coltitle=white,
  title={#1},
  fonttitle=\bfseries,
  boxrule=0.6pt,
  arc=2mm,
  left=12pt,
  right=12pt,
  top=10pt,
  bottom=10pt
}

% ======================
% 易错点框
% ======================

\newtcolorbox{mistakebox}[1][易错点]{
  enhanced,
  breakable,
  colback=ErrorColor!5,
  colframe=ErrorColor,
  colbacktitle=ErrorColor,
  coltitle=white,
  title={#1},
  fonttitle=\bfseries,
  boxrule=0.6pt,
  arc=2mm,
  left=12pt,
  right=12pt,
  top=10pt,
  bottom=10pt
}

% ======================
% 提醒框
% ======================

\newtcolorbox{tipbox}[1][提示]{
  enhanced,
  breakable,
  colback=WarnColor!10,
  colframe=WarnColor,
  colbacktitle=WarnColor,
  coltitle=white,
  title={#1},
  fonttitle=\bfseries,
  boxrule=0.6pt,
  arc=2mm,
  left=12pt,
  right=12pt,
  top=10pt,
  bottom=10pt
}

% ======================
% 章节标题样式
% ======================

\titleformat{\section}
  {\Large\bfseries\color{MainColor}}
  {\thesection}
  {0.8em}
  {}

\titleformat{\subsection}
  {\large\bfseries\color{TitleDark}}
  {\thesubsection}
  {0.8em}
  {}

\titleformat{\subsubsection}
  {\normalsize\bfseries\color{SecondColor!80!black}}
  {\thesubsubsection}
  {0.8em}
  {}

% ======================
% 列表样式
% ======================

\setlist[itemize]{
  leftmargin=2.2em,
  itemsep=0.3em,
  topsep=0.4em
}

\setlist[enumerate]{
  leftmargin=2.2em,
  itemsep=0.3em,
  topsep=0.4em
}

% ======================
% 自定义强调命令
% ======================

\newcommand{\important}[1]{\textbf{\textcolor{MainColor}{#1}}}
\newcommand{\warning}[1]{\textbf{\textcolor{ErrorColor}{#1}}}
\newcommand{\concept}[1]{\textbf{\textcolor{DefineColor}{#1}}}

% ======================
% 正文开始
% ======================
% 一些解释,关于box的说明
% 本文共涉及摘要框summarybox,定义框definitionbox
% 重点框keybox,例题框examplebox,易错点框mistakebox
% 提醒框tipbox 六种框
% ======================


\begin{document}

\topbar{吴恩达机器学习}{学习笔记}

\notetitlecard
  {Study Notes Series}
  {2026 Spring}
  {吴恩达机器学习笔记(1)}
  {}
  {机器学习}
  {C++与草稿纸}

\begin{summarybox}
本章介绍机器学习的基本概念，并以单变量线性回归为例，梳理从建立模型、
定义代价函数到使用梯度下降训练参数的完整过程。
\begin{itemize}[leftmargin=2em, itemsep=2pt, topsep=4pt]
  \item \textbf{机器学习类型：}机器学习主要包括监督学习、非监督学习和强化学习；
  监督学习利用带标签数据学习输入到输出的映射，非监督学习则从无标签数据中发现结构与规律。
  \item \textbf{常见任务：}监督学习可分为预测连续值的回归问题和预测离散类别的分类问题；
  非监督学习的典型任务包括聚类、异常检测和降维。
  \item \textbf{线性回归：}使用 $f_{w,b}(x)=wx+b$ 拟合数据，并通过最小化均方误差代价函数
  $J(w,b)$ 寻找合适的参数 $w$ 和 $b$。
  \item \textbf{梯度下降：}沿 $J(w,b)$ 梯度的反方向迭代更新参数。计算 $w$、$b$ 的偏导数时，
  必须基于同一组更新前的参数，并在得到两个临时值后同时完成更新。
\end{itemize}
\end{summarybox}

\section{机器学习概论}

\begin{definitionbox}[机器学习的定义]
塞缪尔（Arthur Samuel）对机器学习的\textcolor{red}{经典定义}是：

机器学习是使计算机在没有被明确编程的情况下，具备学习能力的研究领域。
\\
Machine Learning is the field of study
that gives computers the ability to learn
without being explicitly programmed.
\end{definitionbox}

\begin{keybox}[机器学习的类型]
如今，机器学习分为三种类型
\begin{itemize}
    \item 监督学习（Supervised Learning）
    \item 非监督学习（Unsupervised Learning）
    \item 强化学习（Reinforcement Learning）
\end{itemize}
其中，监督学习（Supervised Learning）使用较多。
\end{keybox}

\section{监督式机器学习与非监督式机器学习}
\subsection{监督学习}
\begin{definitionbox}[监督学习的关键特征]
给算法一组“输入—正确输出”的样本，让算法学习从输入到输出的映射关系。
\end{definitionbox}


\begin{examplebox}[监督学习的分类和用途]

\begin{center}
    \begin{tabularx}{\linewidth}{|c|c|X|}
        \hline
        类型 & 输出结果 & 典型例子 \\
        \hline
        \textcolor{red}{回归（Regression）} & 连续值 & 预测房价、预测气温、预测成绩 \\
        \hline
        \textcolor{red}{分类（Classification）} & 离散类别 & 判断垃圾邮件、判断肿瘤良恶性、识别数字 \\
        \hline
    \end{tabularx}
\end{center}
\end{examplebox}

\subsection{非监督学习}
\begin{definitionbox}[非监督学习的关键特征]
数据集没有标签，只有特征 $x^{(i)}$。
目标是让算法自行在数据中寻找隐藏的结构、模式或规律。
\end{definitionbox}

\begin{examplebox}[监督学习的分类和主要特征]
\begin{itemize}
  \item \textcolor{red}{聚类(Clustering)}:用于把相同的数据点放在一起
  \item \textcolor{red}{异常检测(Anomaly detection)}:用来寻找特殊数据点
  \item \textcolor{red}{降维(Dimensionality reduction)}:用更少的变量去压缩存储数据

\end{itemize}
\end{examplebox}


\section{回归模型}

\begin{definitionbox}[回归问题]
我们根据之前的数据预测出一个准确的输出值。
\end{definitionbox}
\begin{definitionbox}[分类问题]
当我们想要预测离散的输出值，例如，我们正在寻找癌症肿瘤，并想要确定肿瘤是良性的还是恶性的，这就是0/1离散输出的问题。
\end{definitionbox}
\begin{examplebox}[机器学习的一些符号表示]
\begin{center}
  \begin{tabularx}{\linewidth}{|c|X|}
      \hline
      符号 & 表示意义 \\
      \hline
      $m$ & 训练集中实例的数量\\
      \hline
      $x$ & 特征/输入变量\\
      \hline
      $y$ & 目标变量/输出变量 \\
      \hline
      $(x,y)$ & 训练集中的实例 \\
      \hline
      $(x^{(i)},y^{(i)})$ & 第$i$个观察实例 \\
      \hline
      $h$ & 学习算法的解决方案或函数 \\
      \hline
  \end{tabularx}
\end{center}
\end{examplebox}
\begin{definitionbox}[线性回归]
  用一条直线（或线性函数）（$f(x) = wx + b$）拟合数据，从而预测连续数值。
\end{definitionbox}
\begin{examplebox}[线性回归的代价函数]
  $$Cost \ function J(w,b) = \frac{1}{2m} \sum_{i = 1}^{n} (f_{w,b}^{(i)}(x) - y^{(i)})^2$$
  \[
  Goal : \ \underset{w,b}{minimize} \ J(w,b)
  \]
\end{examplebox}
\section{用梯度下降法训练模型}
\begin{definitionbox}[批量梯度下降公式]
  $$w = w - \alpha \frac{\partial}{\partial w} J(w,b)$$
  其中：$\alpha$表示学习率
\end{definitionbox}
\begin{mistakebox}[参数更新的正确与错误方式]
  \begin{minipage}[t]{0.47\linewidth}
    \vspace{0pt}
    \small
    \textbf{正确：同时更新参数}
    \begin{align*}
      \mathrm{tmp}_w &= w-\alpha\frac{\partial}{\partial w}J(w,b), \\
      \mathrm{tmp}_b &= b-\alpha\frac{\partial}{\partial b}J(w,b), \\
      w &= \mathrm{tmp}_w, \\
      b &= \mathrm{tmp}_b.
    \end{align*}
    两个临时值均由同一组旧参数 $(w,b)$ 计算，之后再同时更新参数。
  \end{minipage}
  \hfill
  \begin{minipage}[t]{0.47\linewidth}
    \vspace{0pt}
    \small
    \textbf{错误：依次更新参数}
    \begin{align*}
      \mathrm{tmp}_w &= w-\alpha\frac{\partial}{\partial w}J(w,b), \\
      w &= \mathrm{tmp}_w, \\
      \mathrm{tmp}_b &= b-\alpha\frac{\partial}{\partial b}J(w,b), \\
      b &= \mathrm{tmp}_b.
    \end{align*}
    计算 $\mathrm{tmp}_b$ 时，$w$ 已是更新后的值，因此两个梯度不在同一参数点上计算。
  \end{minipage}
\end{mistakebox}

\begin{examplebox}[线性回归的梯度下降]
  对于单变量线性回归模型
  \[
    f_{w,b}\bigl(x^{(i)}\bigr)=wx^{(i)}+b,
  \]
  其代价函数是关于 $w$ 和 $b$ 的二元函数：
  \[
    J(w,b)=\frac{1}{2m}\sum_{i=1}^{m}
    \left(f_{w,b}\bigl(x^{(i)}\bigr)-y^{(i)}\right)^2.
  \]
  因此，可将梯度下降写成向量形式：
  \[
    \begin{pmatrix} w \\ b \end{pmatrix}
    \leftarrow
    \begin{pmatrix} w \\ b \end{pmatrix}
    -\alpha\nabla J(w,b),
    \qquad
    \nabla J(w,b)=
    \begin{pmatrix}
      \dfrac{\partial J(w,b)}{\partial w} \\[6pt]
      \dfrac{\partial J(w,b)}{\partial b}
    \end{pmatrix}.
  \]
  对 $J(w,b)$ 分别求偏导，可得批量梯度下降的具体更新公式：
  \begin{align*}
    w &\leftarrow w-\alpha\frac{1}{m}\sum_{i=1}^{m}
    \left(f_{w,b}\bigl(x^{(i)}\bigr)-y^{(i)}\right)x^{(i)}, \\
    b &\leftarrow b-\alpha\frac{1}{m}\sum_{i=1}^{m}
    \left(f_{w,b}\bigl(x^{(i)}\bigr)-y^{(i)}\right).
  \end{align*}
  重复上述过程直至 $J(w,b)$ 收敛。两个偏导数必须基于更新前的同一组
  $(w,b)$ 同时计算。图中的参数与这里的记号对应为
  $\theta_0=b$、$\theta_1=w$。
\end{examplebox}


\end{document}
